%\addbibresource{/home/jorgsk/Dropbox/phdproject/bibtex/jorgsk.bib}
\subsection{Model of initial transcription}
Our model describes the process of initial transcription by considering the
following reactions: i) the NAC ii) backtracking iii) unscrunching and abortive
RNA release (UAR), and iv) promoter escape (Figure
S\ref{fig:model_and_rates}A). We consider backtracking as the first
unscrunching step, also known as backstepping. To find the rate constant of
backtracking we emply a novel method which makes use of abortive probabilities
(APs) calculated from experimental data \cite{hsu_quantitative_1996}. The
method relies on two key points. The first is the model constraint that the
reactions of the NAC and backtracking are in kinetic competition. The second
is the assumption that the probability to backtrack at any given template
position is equal to the AP at that position. Thus we get
\begin{equation*}
    \frac{b_i}{b_i + \text{NAC}_i} = \text{AP}_i,
\end{equation*}
where subscript $i$ indicates position, and NAC and $b$ are the rate constants
the nucleotide addition cycle and backtracking, respectively. From this
expression, we obtain the rate constant for backtracking given a rate constant
of the NAC and an AP:
\begin{equation}
  b_i = \frac{\text{NAC}_i\cdot\text{AP}_i}{1-\text{AP}_i}.
  \label{eq:backtrackingcalc}
\end{equation}
Figure~S\ref{fig:model_and_rates}B shows the initial transcription kinetic
reaction diagram with the rate constants to be estimated, and
Figure~S\ref{fig:backtracking_calc_example} shows example calculations of rate
constants from backtracking.

In this work we run model simulations alternately using APs obtained both for
the absence of GreB (-GreB) and in the presence of GreB (+GreB). When GreB is
present, we note that backtracked complexes may be rescued by GreB stimulated
cleavage of the unaligned 3\ppp end of the transcript. Therefore, APs obtained
in the presence of GreB represent the probability to both backtrack and to
avoid rescue by GreB until abortive RNA release. We therefore assume that
GreB-mediated cleavage and subsequent NACs are rapid steps compared to
unscrunching. This permits using the APs obtained in the presence of GreB as
effective backtracking probabilities. All AP values used in this work are
obtained from Hsu et al. \cite{hsu_initial_2006}.

\subsection{Implementation and rate constant estimation}
A central result in this work is determining the model rate constants. We do
this running multiple simulations with different combinations of rate
constants and fitting to exeprimental data. The data used for rate constant
estimation is the distribution of time spent in abortive cycling on the N25
promoter in the presence of $100\ \mu M$ NTP determined by Revyakin et al.\
\cite{revyakin_abortive_2006}. This data was obtained from 100 individual
initial transcription events. For experiments with single molecules, there is
an inherent stochastic component in the experimental outcome that derives from
the randomness of molecular motion. That this is evident for transcription can
be inferred from the large variation observed in single-molecule experiments
of the speed of transcription elongation \cite{adelman_single_2002,
tolic-norrelykke_diversity_2004}. To account for this randomness, we perform
the kinetic simulations using the Gillespie algorithm for stochastic
simulations of chemical reactions \cite{gillespie_exact_1977}. Specifically,
we make use the version implemented by Maarleveld et al.\ in the StocPY
software \cite{maarleveld_stochpy:_2013}.

The proceedure for rate constant estimation is as follows (rate constants
indicated in Figure S1): First, we assign random values to the three rate
constants NAC ($k_n$), promoter escape ($k_e$) and unscrunching and abortive
RNA release ($k_u$) within certain fixed boundaries (see below). These values
are chosen independently from a uniform distribution. Secondly, we use the
rate constant of the NAC and the APs to calculate the rate constant of
backtracking ($k_b$) using equation \ref{eq:backtrackingcalc}. We then
simulate 100 initial transcription events and calculate the distribution of
time spent in abortive cycling as a result of these specific rate constants.
The result is measured against the empirical distribution
\cite{revyakin_abortive_2006} using the root mean square error (RMSE). By
repeating this procedure several times, we obtain statistics of which
values of rate constants that are associated with the best match with the
experimental data. By virtue of using stochstic simulations, one does not
arrive at a unique set of rate constants that have optimal fit with data; if a
simulation is repeated without changing the rate constants, slightly different
results will be obtained due to the inherent stochastisity of the reaction
process. Therefore, we here find the optimal rate constants by running
multiple simulations and identifying the rate constants that are most often
associated with best fit to data from the distribution of fitness.

The boundaries of the rate constants for parameter estimation were chosen from
extema estimated from experimental data. The speed of the NAC during initial
transcription cannot be much less than 3 nt/s, since Revyakin et al.\ measured
3.5 seconds as the shortest duration of abortive cycling
\cite{revyakin_abortive_2006} to cover the ~ 11 nts required for promoter
escape on N25. At the same time, it should not be larger than at most 25 nt/s,
which is at the upper limit of measured values from single-molecule
experiments \cite{bai_mechanochemical_2007}. For backtracking and abortive RNA
release, it is known that the rate constant must be faster than 1 s$^{-1}$,
this was the time-resolution of the experimental equipment and this reaction
step could not be resolved \cite{revyakin_abortive_2006}. We therefore set the
minimum and maximum values of all rate constants to be 1 s$^{-1}$ and 25
s$^{-1}$, respectively. 

\begin{figure}
	\begin{center}
        \includegraphics{/home/jorgsk/Dropbox/The-Tome/my_papers/kinetic_initiation/illustrations/model_and_rates.pdf}
	\end{center}
    \caption{Kinetic model scheme and reaction rate constants. \textbf{A})
    From the open complex (OC) transcription proceeds by nucleotide additions
    cycles from one initial transcribing complex (IC) to the next, where each
    IC is identified in subscript by the length of its nascent RNA. Initial
    transcription proceeds until the nascent RNA has reached the
    experimentally obtained maximum size of abortive transcript (MSAT)
    \cite{hsu_initial_2006}. For ICs with an RNA of 2 nt in length or longer, there is a
    competition between the nucleotide addition cycle (NAC) and backtracking.
    Backtracking causes termination of transcription, and only further
    backtracking (unscrunching) and abortive RNA release may follow, returning
    RNAP to the open complex. From the open complex forward transcription may
    resume once more. \textbf{B}) Kinetic reactions diagram containing the
    rate constants NAC (denoted here as $k_n$), promoter escape (denoted as
    $k_e$) and unscrunching and abortive release ($k_u$), as well as the rate
    constants of backtracking, $k_{b,i}$.}
    \label{fig:model_and_rates}
\end{figure}
